\section{Multiple Coordinated Views}
\label{sec:theory:mcv}

The Multiple Coordinated Views (MCV) interaction
model~\cite{liu2022} is the principle that connects the
visualisation primitives of Section~\ref{sec:theory:primitives}
and the storytelling principles of
Section~\ref{sec:theory:storytelling} into a coherent
interactive dashboard.

An MCV system is a collection of views over a shared dataset,
with the property that any selection or filter applied in one
view propagates instantaneously to every other view. The
canonical implementation is \emph{brushing-and-linking}: a
brush stroke over one chart highlights the corresponding
regions in all other charts, and the highlight follows the
brush continuously as the analyst moves it. The model
formalises an intuition that has appeared in interactive
analytical systems for several decades; it provides the
vocabulary in which modern dashboard designs can be described
and compared.

\subsection{Implementation requirements}

An MCV system imposes three requirements on its implementation.
First, all views must share a single source of truth for filter
state, so that an update in one view is observable in all the
others without requiring view-to-view notification. Second,
the views must be independently rendered: if one view is
expensive to recompute it must not block the others. Third,
the navigation between views must be one-click rather than
menu-driven, so that analysts can shift their attention from
view to view at the speed of thought.

\subsection{The ControlCenter abstraction}

Liu introduces the \emph{ControlCenter} as the navigation
mechanism for MCV systems with many coordinated views. The
ControlCenter is a persistent UI element (typically an icon
strip or a docked panel) that allows the analyst to switch
between \emph{coordinated view sets} --- groups of views that
share analytical context --- while preserving the underlying
filter state. In Liu's traffic-control deployments the
ControlCenter appears at the left edge of the dashboard with
one icon per view set; clicking an icon swaps the visible
views without disturbing the filter context.

The ControlCenter resolves a long-standing tension in
multi-view dashboards: each analytical question wants its
own view set, but the analyst should not have to re-enter
filters every time they switch from one question to another.
The ControlCenter externalises the question selection from
the views themselves, so the dashboard can offer many
view sets without the cost of redundant filter chrome on each
one.

\subsection{Brushing-and-linking versus selection sharing}

Two flavours of MCV propagation appear in the literature.
\emph{Brushing-and-linking} highlights the same data points
across all views simultaneously: hovering a point on a scatter
plot illuminates the corresponding row in a table and the
corresponding region on a map. \emph{Selection sharing}
propagates a discrete selection event: clicking a tile in a
treemap commits a new value to the shared filter state, which
all views then re-render against. Brushing favours hover-driven
exploration; selection sharing favours click-driven drill-down.
A practical dashboard typically supports both, with brushing
on hover and selection sharing on click.

The dashboard implements MCV at full scope through what
Section~\ref{sec:proposed:invariants} denotes A1 linked brushing.
Within a single card, selecting an entity (a route row in a
heatmap, a bubble in a quadrant, a hex cell in a map) propagates
the selection to every other chart on the same card. The
implementation is server-side: a click writes to
\texttt{st.session\_state[\textquotesingle highlight\textquotesingle]},
\texttt{st.rerun()} triggers a re-render of all charts on the
page, and each chart reads the highlight on render. Cross-card
linking is reserved for explicit drill-down deep-links rather
than implicit propagation, to keep the per-card mental model
local.
